\section*{Bab \ref{ch:g9:sqrt} --- Akar Kuadrat}

\begin{solution}{exo:g9:sqrt:1}
$\sqrt{64} = 8$; \quad $\sqrt{100} = 10$; \quad
$\sqrt{0.09} = 0.3$ (karena $0.3^2 = 0.09$); \quad
$\sqrt{\frac{1}{25}} = \frac15$; \quad
$\left(\sqrt{13}\right)^2 = 13$; \quad
$\sqrt{(-4)^2} = \sqrt{16} = 4 = \abs{-4}$.
\end{solution}

\begin{solution}{exo:g9:sqrt:2}
$\sqrt{50} = \sqrt{25 \times 2} = 5\sqrt2$; \quad
$\sqrt{45} = \sqrt{9 \times 5} = 3\sqrt5$; \quad
$\sqrt{98} = \sqrt{49 \times 2} = 7\sqrt2$; \quad
$\sqrt{300} = \sqrt{100 \times 3} = 10\sqrt3$.
\end{solution}

\begin{solution}{exo:g9:sqrt:3}
$\sqrt2 \times \sqrt{18} = \sqrt{36} = 6$.

$\dfrac{\sqrt{75}}{\sqrt3} = \sqrt{\dfrac{75}{3}} = \sqrt{25} = 5$.

$\sqrt5 \times \sqrt{20} = \sqrt{100} = 10$.

$\dfrac{\sqrt8}{\sqrt{50}} = \sqrt{\dfrac{8}{50}} = \sqrt{\dfrac{4}{25}}
= \dfrac25$.
\end{solution}

\begin{solution}{exo:g9:sqrt:4}
$x^2 = 36$: $x = 6$ atau $x = -6$.

$x^2 = 11$: $x = \sqrt{11}$ atau $x = -\sqrt{11}$.

$x^2 + 4 = 0$ berarti $x^2 = -4 < 0$: tidak ada penyelesaiannya.

$2x^2 = 50$: $x^2 = 25$, jadi $x = 5$ atau $x = -5$.
\end{solution}

\begin{solution}{exo:g9:sqrt:5}
$\sqrt{20} + \sqrt{45} = 2\sqrt5 + 3\sqrt5 = 5\sqrt5$.

$3\sqrt8 - \sqrt{18} + \sqrt2 = 3 \times 2\sqrt2 - 3\sqrt2 + \sqrt2
= (6 - 3 + 1)\sqrt2 = 4\sqrt2$.
\end{solution}

\begin{solution}{exo:g9:sqrt:6}
$\left(\sqrt3 + 1\right)^2 = 3 + 2\sqrt3 + 1 = 4 + 2\sqrt3$.

$\left(2\sqrt5 - 3\right)^2 = 4 \times 5 - 2 \times 3 \times 2\sqrt5 + 9
= 29 - 12\sqrt5$.

$\left(\sqrt6 + \sqrt2\right)\left(\sqrt6 - \sqrt2\right) = 6 - 2 = 4$.
\end{solution}

\begin{solution}{exo:g9:sqrt:7}
Sisinya: $\sqrt{6400} = 80$ m. Diagonalnya (dengan Pythagoras):
$\sqrt{80^2 + 80^2} = 80\sqrt2 \approx 113$ m.
\end{solution}

\begin{solution}{exo:g9:sqrt:8}
$\dfrac{1}{\sqrt2} = \dfrac{1 \times \sqrt2}{\sqrt2 \times \sqrt2}
= \dfrac{\sqrt2}{2}$.

$\dfrac{6}{\sqrt3} = \dfrac{6\sqrt3}{3} = 2\sqrt3$.
\quad
$\dfrac{10}{\sqrt5} = \dfrac{10\sqrt5}{5} = 2\sqrt5$.
\end{solution}

\begin{solution}{exo:g9:sqrt:9}
\emph{1. Benar}: inilah aturan \omterm{def:g2:mult:def}{hasil kalinya} (\cref{thm:g9:sqrt:rules}).

\emph{2. Salah}: $\sqrt{9 + 16} = 5$ tetapi $\sqrt9 + \sqrt{16} = 7$.

\emph{3. Salah} untuk $x$ yang negatif: $\sqrt{(-4)^2} = 4 \neq -4$.
Pernyataan yang benar adalah $\sqrt{x^2} = \abs{x}$.

\emph{4. Benar}: $\left(3\sqrt2\right)^2 = 9 \times 2 = 18$.
\end{solution}

\begin{solution}{exo:g9:sqrt:10}
$\left(2 + \sqrt3\right)^2 = 4 + 4\sqrt3 + 3 = 7 + 4\sqrt3$. Jadi
$x = \sqrt{7 + 4\sqrt3} = \sqrt{\left(2+\sqrt3\right)^2} = 2 + \sqrt3$
(yaitu bilangan taknegatif, jadi akarnya tinggal membuang kuadratnya).

Serupa dengan itu $\left(\sqrt5 - 2\right)^2 = 5 - 4\sqrt5 + 4 = 9 -
4\sqrt5$, dan $\sqrt5 - 2 > 0$, jadi
$\sqrt{9 - 4\sqrt5} = \sqrt5 - 2$ (bukan $2 - \sqrt5$, yang bernilai
negatif!).
\end{solution}

\begin{solution}{pb:g9:sqrt:1}
\textbf{1.} $1.4^2 = 1.96 < 2 < 2.25 = 1.5^2$;
$1.41^2 = 1.9881 < 2 < 2.0164 = 1.42^2$. Tiga desimal:
$1.414^2 = 1.999396 < 2 < 2.002225 = 1.415^2$, jadi
$1.414 < \sqrt2 < 1.415$.

\textbf{2.} Mengkuadratkan $\sqrt2 = \frac ab$ memberi
$2 = \frac{a^2}{b^2}$, karena itu $a^2 = 2b^2$: jadi bilangan $a^2$
adalah dua kali sebuah bilangan bulat --- yaitu \emph{\omterm{def:g3:numbers:evenodd}{genap}}.

\textbf{3.} Kalau $a$ \omterm{def:g3:numbers:evenodd}{ganjil}, maka $a^2$ akan \omterm{def:g3:numbers:evenodd}{ganjil}
(\cref{pb:g8:pythagoras:1}); karena $a^2$ \omterm{def:g3:numbers:evenodd}{genap}, maka $a$ \omterm{def:g3:numbers:evenodd}{genap}:
$a = 2k$. Dengan menyulihkannya: $4k^2 = 2b^2$, jadi $b^2 = 2k^2$
\omterm{def:g3:numbers:evenodd}{genap}, dan menurut kenyataan \omterm{def:g3:numbers:evenodd}{ganjil-genap} yang sama $b$ pun \omterm{def:g3:numbers:evenodd}{genap}.

\textbf{4.} $a$ dan $b$ yang keduanya \omterm{def:g3:numbers:evenodd}{genap} berarti keduanya \omterm{def:g6:wholes:divisible}{habis
dibagi} $2$ --- tetapi $\frac ab$ tadi paling sederhana, sehingga tidak
berbagi \omterm{def:g9:arith:divisor}{pembagi} bersama. Kontradiksi: \omterm{def:g9:fractions:fraction}{pecahan} yang diandaikan itu
tidak dapat ada. \emph{Teorema: $\sqrt2$ irasional.}

\textbf{5.} Kontradiksinya sepenuhnya hidup pada kata ``paling
sederhana'': tanpa kata itu, ``$a$ dan $b$ keduanya \omterm{def:g3:numbers:evenodd}{genap}'' tidak
bertentangan dengan apa pun. Mengandaikan bentuk paling sederhana itu
sah karena setiap \omterm{def:g9:fractions:fraction}{pecahan} \emph{punya} bentuk paling sederhana
(\cref{met:g9:arith:simplify}: bagilah dengan FPB-nya) --- kalau
$\sqrt2$ berupa \omterm{def:g9:fractions:fraction}{pecahan} apa pun, ia akan berupa \omterm{def:g9:fractions:fraction}{pecahan} yang paling
sederhana, dan \omterm{def:g9:fractions:fraction}{pecahan} seperti itu mustahil.

\textbf{6.} Pada pemfaktoran prima, kuadrat membawa \omterm{def:g8:powers:def}{eksponen} yang
\omterm{def:g3:numbers:evenodd}{genap} (\cref{pb:g9:arith:1}). Pada $a^2 = 3b^2$, \omterm{def:g8:powers:def}{eksponen}
$3$ \omterm{def:g3:numbers:evenodd}{genap} di ruas kirinya, tetapi \omterm{def:g3:numbers:evenodd}{ganjil} di ruas kanannya (\omterm{def:g3:numbers:evenodd}{genap} dari
$b^2$, ditambah satu). Tidak ada bilangan yang punya dua pemfaktoran
(\cref{thm:g9:arith:factorization}): jadi kontradiksi, dan
$\sqrt3$ irasional.

\textbf{7.} Pada $a^2 = n b^2$, alasannya menemukan \omterm{def:g9:arith:prime}{bilangan prima}
yang sifat \omterm{def:g3:numbers:evenodd}{ganjil-genap} \omterm{def:g8:powers:def}{eksponennya} bertentangan persis ketika ada
\omterm{def:g9:arith:prime}{bilangan prima} yang muncul pada $n$ dengan \omterm{def:g8:powers:def}{eksponen} yang
\emph{\omterm{def:g3:numbers:evenodd}{ganjil}}. Kalau setiap \omterm{def:g8:powers:def}{eksponen} $n$ \omterm{def:g3:numbers:evenodd}{genap}, maka $n$ adalah
kuadrat sempurna ($n = m^2$, dan $\sqrt n = m$ bilangan bulat). Hasil
lengkapnya: $\sqrt n$ irasional untuk setiap bilangan bulat $n$ yang
bukan kuadrat sempurna ---
$\sqrt2, \sqrt3, \sqrt5, \sqrt6, \sqrt7, \sqrt8, \sqrt{10},
\dots$

\textbf{8.} Dari $r x = q$ dengan $r \neq 0$:
$x = \frac qr$, yaitu \omterm{def:g3:division:remainder}{hasil bagi} bilangan rasional, jadi rasional.
Jadi kalau $x$ irasional, $r x$ tidak dapat rasional: $2\sqrt2$ dan
$\frac{\sqrt2}{2}$ irasional (dengan pengali $2$ dan
$\frac12$).

\textbf{9.} Kalau $1 + \sqrt2 = q$ rasional, maka
$\sqrt2 = q - 1$ akan rasional: kontradiksi. Kalau
$s = \sqrt2 + \sqrt3$ rasional, maka
$s^2 = 2 + 2\sqrt6 + 3 = 5 + 2\sqrt6$ akan rasional juga, sehingga
memberi $\sqrt6 = \frac{s^2 - 5}{2}$ yang rasional --- tetapi
$\sqrt6$ irasional (pertanyaan 7). Jadi $\sqrt2 + \sqrt3$ irasional.

\textbf{10.} \omterm{def:g1:addition:def}{Jumlahnya}: $\sqrt2$ dan $-\sqrt2$ (atau $\sqrt2$ dan
$1 - \sqrt2$) masing-masing irasional, dengan \omterm{def:g1:addition:def}{jumlah} rasional $0$ dan
$1$. \omterm{def:g2:mult:def}{Hasil kalinya}: $\sqrt2 \times \sqrt2 = 2$. Keirasionalan dapat
menguap dalam aritmetika --- karena itulah buktinya satu per satu.

\textbf{11.} Dari $x = 1$: rata-rata $1$ dan $\frac21 = 2$:
$x_1 = \frac32$. Lalu rata-rata $\frac32$ dan
$\frac{2}{3/2} = \frac43$:
$x_2 = \frac12\left(\frac32 + \frac43\right)
= \frac12 \cdot \frac{17}{6} = \frac{17}{12} \approx 1.4167$.

\textbf{12.} $x_3 = \frac12\left(\frac{17}{12} +
\frac{24}{17}\right) = \frac12 \cdot
\frac{289 + 288}{204} = \frac{577}{408} = 1.4142156\ldots$
Dibandingkan $\sqrt2 = 1.4142135\ldots$: lima desimalnya benar sesudah
tiga langkah --- resepnya kira-kira melipatduakan desimal yang benar
pada tiap putarannya.

\textbf{13.} Kalau $x^2 > 2$ (tebakannya terlalu besar), maka
$x > \sqrt2$, dan $\frac2x < \frac{2}{\sqrt2} = \sqrt2$: jadi sisi
pasangannya terlalu kecil. Kalau $x^2 < 2$, ketaksamaannya terbalik.
Bagaimanapun juga $\sqrt2$ duduk di antara $x$ dan $\frac2x$, jadi
rata-ratanya --- yaitu tebakan berikutnya --- lebih dekat daripada
sisi yang lebih buruk dari keduanya: persegi panjang berluas $2$ itu
makin mempersegi pada tiap langkahnya.

\textbf{14.} $9 - 8 = 1$; \ $289 - 288 = 1$; \
$332\,929 - 332\,928 = 1$. Setiap tebakan Heron $\frac ab$
memenuhi $a^2 - 2b^2 = 1$: Bagian I membuktikan bahwa mengenai $0$
itu mustahil, dan \omterm{def:g9:fractions:fraction}{pecahan} ini meleset dari yang mustahil tepat
satu satuan --- sedekat yang pernah dapat dicapai bilangan bulat.

\textbf{15.} Garis bilangannya pasti memuat bilangan di luar
\omterm{def:g9:fractions:fraction}{pecahan} --- yaitu panjang seperti $\sqrt2$, yang penulisan desimalnya
berjalan selamanya tanpa berulang: itulah \emph{bilangan real}.
Pembangunannya yang jujur adalah kisah untuk buku Sekolah Menengah
dan, dengan kekakuan penuh, buku universitasnya.

\end{solution}
